% =============================================================
%  CAPÍTULO I — MARCO CONTEXTUAL DEL PROYECTO TECNOLÓGICO
% =============================================================
% =============================================================
%  CAPÍTULO I — MARCO CONTEXTUAL DEL PROYECTO TECNOLÓGICO
% =============================================================

\section{Problema de investigación}

Las transformaciones educativas de las últimas décadas han desplazado el énfasis tradicional centrado en la transmisión de contenidos hacia un paradigma orientado al desarrollo de competencias, entendidas como la capacidad del estudiante para movilizar conocimientos, habilidades y actitudes en situaciones auténticas \cite{black1998inside,unesco2022contract,oecd2019compass}. Este giro paradigmático ha sido impulsado simultáneamente por organismos internacionales como la OCDE, la UNESCO y el Banco Mundial, y por las reformas curriculares regionales que reconocen en la evaluación por competencias un instrumento estratégico para elevar la calidad educativa \cite{voinea2025competency,velezvelez2024curriculo}. Sin embargo, la traducción operativa de este enfoque al aula sigue siendo un reto abierto, particularmente en la educación básica elemental, donde concurren factores cognitivos, emocionales y de accesibilidad que demandan instrumentos especialmente diseñados.

\subsection{Planteamiento del problema}

Es innegable la creciente importancia de implementar evaluaciones por competencias curriculares en el contexto educativo. En un mundo cada vez más orientado hacia el desarrollo de habilidades, estas evaluaciones se posicionan como elementos fundamentales para medir con precisión las destrezas de los estudiantes en relación con los objetivos y contenidos del plan de estudios \cite{black1998inside,maskos2025fa}.

Las evaluaciones por competencias se presentan como herramientas esenciales que van más allá de evaluar conocimientos teóricos, abordando la capacidad práctica y la aplicación efectiva de conceptos. Este enfoque integral no solo proporciona una evaluación más precisa del rendimiento estudiantil, sino que también contribuye a preparar a los estudiantes con las habilidades necesarias para enfrentar los retos del entorno educativo actual y futuro \cite{voinea2025competency,oecd2019compass}.

La implementación efectiva y generalizada de las evaluaciones por competencia presenta desafíos significativos. Aunque la tecnología ha avanzado, la integración de estas evaluaciones en las prácticas educativas sigue siendo limitada \cite{velezvelez2024curriculo,niss2019math}. La falta de plataformas especializadas y adaptadas a las necesidades específicas de la evaluación por competencias dificulta su aplicación amplia y eficiente. Asimismo, la diversidad de áreas curriculares y la variabilidad en los estilos de aprendizaje de los estudiantes plantean la necesidad de una solución integral que pueda adaptarse a distintos contextos educativos. La escasez de herramientas que permitan la aplicación de evaluaciones por competencias de manera personalizada y detallada, así como la falta de recursos que faciliten la interpretación y el análisis de los resultados por parte de los docentes, son obstáculos adicionales \cite{gonzalez2021ai,maksimchuk2025genai}.

\subsubsection{Diagnóstico}

Actualmente, el proceso de evaluación en la educación básica elemental, específicamente en el área de matemáticas, se ve limitado por métodos tradicionales que priorizan la memorización sobre el desempeño. Se observa que los docentes carecen de herramientas tecnológicas diseñadas específicamente para medir competencias curriculares de forma individualizada \cite{velezvelez2024curriculo,evans2020cbe}, lo que genera una brecha entre lo que el estudiante ``sabe'' teóricamente y lo que puede ``hacer'' en contextos prácticos.

A esta limitación instrumental se suman barreras cognitivas y emocionales propias de este nivel formativo. La literatura especializada documenta que la ansiedad matemática constituye un factor de alta prevalencia en estudiantes de educación básica, capaz de reducir la capacidad de la memoria de trabajo y deteriorar significativamente el desempeño en tareas aritméticas \cite{passolunghi2016math,ramirez2018mathanxiety,mitchell2022prevalence}. De manera complementaria, las dificultades de lectoescritura propias de los primeros años escolares condicionan la comprensión autónoma de los enunciados, situación que la literatura sobre accesibilidad cognitiva propone mitigar mediante apoyos visuales como pictogramas, los cuales reducen la carga cognitiva y facilitan la asociación entre representaciones simbólicas y elementos del mundo real \cite{mayer2014multimedia,guillomia2021cognitive}.

Desde la perspectiva del docente, el panorama no es menos crítico. La gestión manual de evaluaciones, calificaciones y reportes individualizados resulta operativamente inviable cuando se atiende a un número considerable de estudiantes, lo que limita la posibilidad de proveer retroalimentación oportuna y orientada al progreso, siendo esta uno de los factores con mayor evidencia empírica de impacto positivo en el aprendizaje \cite{hattie2007feedback,maskos2025fa}. Además, la ausencia de plataformas que centralicen la información del desempeño impide la construcción de una analítica de aprendizaje útil para la toma de decisiones pedagógicas, dejando al profesorado sin insumos confiables para detectar oportunamente vacíos de aprendizaje y planificar acciones de refuerzo \cite{gonzalez2021ai,maksimchuk2025genai}.

\subsubsection{Pronóstico}

De no implementarse soluciones tecnológicas que faciliten la evaluación por competencias, persistirá la desconexión entre el currículo académico y las habilidades reales de los estudiantes. Esto resultará en un bajo rendimiento en pruebas estandarizadas y una desmotivación creciente en el aprendizaje de las matemáticas. Los resultados del Programme for International Student Assessment (PISA) 2022 ya advierten sobre esta tendencia: en América Latina y el Caribe, tres de cada cuatro estudiantes no alcanzan el nivel mínimo de competencia matemática, frente al 31\,\% promedio de la OCDE \cite{oecd2023pisa,bid2023pisa,unesco2024lac}. La continuidad de prácticas evaluativas tradicionales, sin instrumentos especializados para la educación básica elemental, profundizará esta brecha y trasladará al nivel secundario y superior estudiantes con vacíos no remediados en las competencias matemáticas fundamentales.

A largo plazo, la falta de una medición precisa de las destrezas individuales fomentará la exclusión educativa por una doble vía. Por un lado, el sistema no podrá adaptarse a los diferentes ritmos y estilos de aprendizaje, perpetuando la desventaja de aquellos estudiantes que requieren apoyos diferenciados o presentan barreras de lectoescritura, ansiedad matemática o necesidades educativas especiales \cite{passolunghi2016math,guillomia2021cognitive}. Por otro lado, la sobrecarga operativa del docente y la ausencia de analítica integrada limitarán su capacidad de identificar tempranamente a los estudiantes en riesgo, debilitando las posibilidades de intervención pedagógica oportuna \cite{hattie2007feedback,maskos2025fa}. Adicionalmente, la persistencia de esta brecha tecnológica entre los procesos educativos formales y las herramientas que sí utilizan los estudiantes en contextos no escolares profundizará la desconexión entre la experiencia digital cotidiana del niño y la institución educativa, limitando su preparación para los retos académicos y profesionales del futuro \cite{voinea2025competency,unesco2022contract}.

\subsection{Formulación del problema}

¿Cómo superar los desafíos relacionados con la evaluación de competencias curriculares en matemáticas, asegurando precisión, inclusión y adaptabilidad a las necesidades individuales de los estudiantes de educación básica elemental?

\subsection{Sistematización del problema}

\begin{itemize}[leftmargin=*]
  \item ¿Qué necesidades individuales y criterios deben considerarse en el diseño de una herramienta que permita una medición precisa y eficaz de las competencias de los estudiantes?
  \item ¿Qué elementos utilizar en la formulación de las preguntas para el proceso evaluativo que asegure una evaluación integral de las competencias de los estudiantes?
  \item ¿Qué factores deben considerarse en los procesos de evaluación por competencias para garantizar su efectividad en contextos educativos reales?
\end{itemize}

\section{Objetivos}
\label{sec:objetivos}
Los objetivos del presente proyecto tecnológico se formulan como respuesta directa al problema de investigación planteado y a las preguntas de sistematización derivadas. Se estructuran en torno a tres etapas secuenciales y complementarias: una primera fase de levantamiento de información orientada a comprender en profundidad las necesidades de los actores educativos involucrados; una segunda fase de diseño y desarrollo, en la cual se materializa una solución tecnológica inclusiva fundamentada en dichos hallazgos; y una tercera fase de validación empírica en entornos reales, que permite contrastar la efectividad de la propuesta y orientar futuras mejoras. Esta articulación garantiza la trazabilidad entre el problema diagnosticado, la solución desarrollada y la evidencia recopilada para su validación.

\subsection{Objetivo General}

Desarrollar una aplicación web para superar los desafíos relacionados con la evaluación de competencias curriculares en matemáticas, asegurando precisión, inclusión y adaptabilidad a las necesidades individuales de los estudiantes de educación básica elemental aplicando el uso de pictogramas.

\subsection{Objetivos Específicos}

\begin{itemize}[leftmargin=*]
  \item Identificar las necesidades de los estudiantes y las limitaciones que enfrentan los docentes en la evaluación por competencias, mediante entrevistas a profesionales educativos, con el fin de establecer los requisitos funcionales y pedagógicos para el diseño de la aplicación.
  \item Diseñar la estructura y funcionalidades de una aplicación web que integre múltiples tipos de preguntas, ajustes de dificultad basados en el currículo, validación de respuestas y elementos de accesibilidad como pictogramas, con el fin de atender las necesidades individuales de los estudiantes en la evaluación por competencias.
  \item Validar la funcionalidad de la aplicación en entornos educativos reales, observando su desempeño durante la administración de evaluaciones por competencias y recopilando percepciones de docentes y estudiantes para orientar mejoras que fortalezcan su aplicación en el contexto escolar.
\end{itemize}

\section{Justificación}

La evaluación por competencias se ha posicionado como un pilar fundamental en la transformación educativa contemporánea, desplazando el interés desde la mera memorización de contenidos hacia la capacidad de movilizar conocimientos en situaciones prácticas \cite{black1998inside,oecd2019compass,voinea2025competency}. Sin embargo, en el contexto de la educación básica ecuatoriana, existe una brecha significativa entre estos postulados teóricos y la realidad del aula \cite{velezvelez2024curriculo,niss2019math,velezvelez2024curriculo}. Las herramientas actuales, como SCALA o CAT, están diseñadas predominantemente para la educación superior, dejando desatendido el nivel básico elemental \cite{gonzalez2021ai}. Este vacío instrumental impide que las instituciones educativas locales implementen modelos evaluativos modernos, perpetuando prácticas tradicionales que no logran medir con fidelidad el desarrollo de destrezas críticas en asignaturas fundamentales como las matemáticas.

Desde una perspectiva práctica, el desarrollo de este proyecto responde a la necesidad urgente de dotar a los docentes de un instrumento funcional y especializado que automatice y optimice el proceso evaluativo. La utilidad de la aplicación propuesta radica en su capacidad para ofrecer reportes detallados y personalizados sobre el desempeño del estudiante, algo que manualmente resulta inviable para un profesor con numerosos alumnos \cite{hattie2007feedback,maksimchuk2025genai}. La integración de retroalimentación generada por inteligencia artificial constituye, además, uno de los factores con mayor evidencia empírica de impacto positivo en el aprendizaje, según los meta-análisis disponibles \cite{hattie2007feedback,gonzalez2021ai}.

Por esta razón, se plantea la creación de una aplicación web inclusiva destinada a la realización de evaluaciones por competencias curriculares a estudiantes de primaria, centrada en la disciplina de Matemáticas. La propuesta incluye la segmentación de las evaluaciones por módulos, lo que permite adaptar la dificultad de acuerdo con la estructura de los cursos, así como la integración de una variedad de tipos de preguntas para captar el interés y la participación activa del estudiante \cite{debrenti2024gbl,sayed2023adaptive}. Adicionalmente, la incorporación de pictogramas como soporte visual responde a la evidencia que demuestra su capacidad para reducir la carga cognitiva, facilitar la comprensión autónoma y disminuir la ansiedad asociada a las tareas matemáticas en niños \cite{mayer2014multimedia,guillomia2021cognitive,passolunghi2016math}. Desde la perspectiva del docente, se subraya la generación de informes individuales detallados para que puedan interpretar y aplicar los resultados en el proceso educativo \cite{maskos2025fa}.

La rigurosidad del proceso de desarrollo se garantiza mediante la adopción de la metodología ágil Extreme Programming (XP). A diferencia de los modelos en cascada o metodologías orientadas al hardware, XP es idónea para este contexto educativo debido a su enfoque en la simplicidad, la retroalimentación constante y la entrega de valor iterativa \cite{beck2004xp,dyba2008agile,shrivastava2021xpreview}. Dado que las necesidades pedagógicas pueden ser dinámicas, XP permite validar continuamente las funcionalidades --como los tipos de preguntas y la visualización de reportes-- directamente con los usuarios finales, asegurando que el software resultante no sea solo un ejercicio de programación, sino una solución validada que responde fielmente a los problemas reales detectados en la fase de diagnóstico.

Adicionalmente, este proyecto se inscribe en el marco de los Objetivos de Desarrollo Sostenible de la Agenda 2030 de la Organización de las Naciones Unidas, particularmente en el ODS~4 ``Educación de Calidad'', cuya meta 4.1 promueve que todos los niños y niñas terminen la enseñanza primaria y secundaria con resultados de aprendizaje pertinentes y efectivos, y cuya meta 4.5 demanda eliminar las disparidades de género y garantizar el acceso igualitario a la educación para las personas en situación de vulnerabilidad \cite{onu2015ods}. En el ámbito nacional, la propuesta se alinea con el Plan Nacional de Desarrollo ``Creando Oportunidades'' 2024-2025, específicamente con el Eje Social y su Objetivo~5, orientado a proteger a las familias y garantizar una educación de calidad con énfasis en el cierre de brechas \cite{senplades2024pnd}. De igual manera, el proyecto se articula con el Currículo Priorizado emitido por el Ministerio de Educación del Ecuador, el cual estructura el área de Matemáticas en torno a destrezas con criterio de desempeño que requieren instrumentos evaluativos especializados para su correcta medición \cite{mineduc2021curriculo}. La articulación de Kiddy Quiz con estos tres marcos~---~internacional, nacional y curricular~---~refuerza la pertinencia social del proyecto y su contribución concreta a la modernización del sistema educativo ecuatoriano.

Desde una perspectiva social, el desarrollo de Kiddy Quiz trasciende la
creación de una herramienta tecnológica convencional al constituirse como
un mecanismo de equidad orientado a romper las barreras de aprendizaje
imperantes en el modelo de evaluación tradicional. El proyecto se
fundamenta en la necesidad de transformar la tecnología en un puente
hacia la inclusión, proporcionando un entorno digital diseñado
específicamente para atender la diversidad cognitiva del estudiantado
de educación básica. Mediante una combinación estructurada de apoyos
visuales ---basados en pictogramas--- y estímulos auditivos a través de
síntesis de voz, la plataforma garantiza que la evaluación de
competencias matemáticas sea un proceso accesible y equitativo. Este
enfoque permite que el éxito académico no esté condicionado por las
habilidades de lectoescritura ni por el ritmo de procesamiento de
información del niño, dando cumplimiento así al derecho fundamental de
acceso a una educación de calidad sin exclusiones
\cite{unesco2017inclusion,mineduc2023dua}.